\documentclass[12pt,a4paper]{article}
\usepackage{ugmath}
\usepackage{float}
\usepackage{placeins}
\usepackage [spanish]{babel}
\usepackage[labelfont=bf,labelsep=space]{caption}
\newcommand{\paren}[1]{\left( #1 \right)}
\newcommand{\alumno}{Ricardo León Martínez}
\newcommand{\materia}{Calculo Diferencial e Integral II}
\newcommand{\profesor}{Fernando Nuñez Medina}
\newcommand{\tarea}{Tarea 11}
\newcommand{\fecha}{9/5/2026}

\begin{document}
    \begin{center}
    {\large\textbf{UNIVERSIDAD DE GUANAJUATO}}\\[0.3cm]
    {\normalsize\textbf{DIVISIÓN DE CIENCIAS NATURALES Y EXACTAS}}\\
    {\normalsize\textbf{CAMPUS GUANAJUATO}}\\[1cm]

    {\Large\textbf{\tarea\ (\materia)}}\\[1cm]
    \end{center}

    \textbf{Nombre:} \alumno \hfill 
    \textbf{Fecha:} \fecha \hfill 
    \textbf{Calificación:} \rule{3cm}{0.4pt} \\[0.3cm]

    \begin{mdframed}[style=mdpinkbox,frametitle={Ejercicio 1}]
        Prueba el criterio de los límites superiores e inferiores (proposición 57).
    \end{mdframed}

    \begin{proof}
        Definamos
        \begin{equation*}
            s_{n}:=\sup_{k\geq n}a_{k},\quad i_{n}:=\inf_{k\geq n}a_{k}.
        \end{equation*}
        Entonces $(s_{n})$ es decreciente, $(i_{n})$ es creciente y para todo $n$,
        \begin{equation*}
            i_{n}\leq a_{n}\leq s_{n}.
        \end{equation*}
        Supongamos que $a_{n}\to L$. Sea $\varepsilon\gt0$. Como $a_{n}\to L$, existe $N\in\N$ tal que
        \begin{equation*}
            n\geq N\implies\abs{a_{n}-L}\lt\varepsilon.
        \end{equation*}
        Equivalentemente,
        \begin{equation*}
            L-\varepsilon\lt a_{n}\lt L+\varepsilon.
        \end{equation*}
        Fijemos $n\geq N$. Entonces para todo $k\geq n$,
        \begin{equation*}
            L-\varepsilon\lt a_{k}\lt L+\varepsilon.
        \end{equation*}
        Por propiedades de supremo e infimo,
        \begin{equation*}
            \inf_{k\geq n}a_{k}\geq L-\varepsilon,\quad\text{ y }\quad\sup_{k\geq n}a_{k}\leq L+\varepsilon.
        \end{equation*}
        Ademas, como todos los $a_{k}$ pertenecen al intervalo abierto anterior,
        \begin{equation*}
            \inf_{k\geq n}a_{k}\leq L+\varepsilon,\quad sup_{k\geq n}a_{k}\geq L-\varepsilon.
        \end{equation*}
        Por lo tanto,
        \begin{equation*}
            L-\varepsilon\leq\liminf a_{n}\leq\limsup a_{n}\leq L+\varepsilon.
        \end{equation*}
        como esto vale para todo $\varepsilon\gt0$, concluimos que
        \begin{equation*}
            \liminf a_{n}=L,\qquad\limsup a_{n}=L.
        \end{equation*}
        Supongamos ahora que
        \begin{equation*}
            \limsup a_{n}=\liminf a_{n}=L.
        \end{equation*}
        Es decir,
        \begin{equation*}
            s_{n}\to L,\quad i_{n}\to L.
        \end{equation*}
        Como 
        \begin{equation*}
            i_{n}\leq a_{n}\leq s_{n},
        \end{equation*}
        y ademas $i_{n}\to L$, $s_{n}\to L$, por el lema del sandwich se sigue inmediatamente que
        \begin{equation*}
            a_{n}\to L.
        \end{equation*}
        Concluimos que
        \begin{equation*}
            \lim_[n\to\infty]a_{n}=L\text{ ssi }\limsup_{n\to\infty}a_{n}=\liminf_{n\to\infty}a_{n}=L.
        \end{equation*}
    \end{proof}

    \begin{mdframed}[style=mdpinkbox,frametitle={Ejercicio 2}]
        Muestra un ejemplo de un agrupamiento convergente de una serie divergente.
    \end{mdframed}

    \begin{mdframed}[style=mdpinkbox,frametitle={Ejercicio 3}]
        Muestra dos series convergentes $\sum_{n=1}^{\infty}a_{n}$ y $\sum_{n=1}^{\infty}b_{n}$ tales que
        \begin{equation*}
            \sum_{n=1}^{\infty}a_{n}b_{n}\neq\paren{\sum_{n=1}^{\infty}a_{n}}\paren{\sum_{n=1}^{b_{n}}}.
        \end{equation*}
    \end{mdframed}

    \begin{solucion}
        Tomemos
        \begin{equation*}
            a_{n}=b_{n}=\frac{1}{2^{n}}.
        \end{equation*}
        Entonces
        \begin{equation*}
            \sum_{n=1}^{\infty}a_{n}=\sum_{n=1}^{\infty}b_{n}=\sum_{n=1}^{\infty}\frac{1}{2^{n}}=1.
        \end{equation*}
        Por tanto,
        \begin{equation*}
            \paren{\sum_{n=1}^{\infty}a_{n}}\paren{\sum_{n=1}^{\infty}b_{n}}=1.
        \end{equation*}
        Ahora,
        \begin{equation*}
            a_{n}b_{n}=\frac{1}{4^{n}},
        \end{equation*}
        asi que
        \begin{equation*}
            \sum_{n=1}^{\infty}a_{n}b_{n}=\sum_{n=1}^{\infty}\frac{1}{4^{n}}=\frac{\frac{1}{4}}{1-\frac{1}{4}}=\frac{1}{3}.
        \end{equation*}
        Luego,
        \begin{equation*}
            \sum_{n=1}^{\infty}a_{n}b_{n}=\frac{1}{3}\neq1=\paren{\sum_{n=1}^{\infty}a_{n}}\paren{\sum_{n=1}^{\infty}b_{n}}.
        \end{equation*}
    \end{solucion}

    \begin{mdframed}[style=mdpinkbox,frametitle={Ejercicio 4}]
        Determina la convergencia de las series siguientes:
        \begin{multicols}{2}
            \begin{enumerate}
                \item[(a)] $\sum_{n=0}^{\infty}\frac{c^{n}}{n!}$ ($c$ es constante).
                \item[(b)] $\sum_{n=1}^{\infty}\frac{n}{n^{2}+1}$.
                \item[(c)] $\sum_{n=2}^{\infty}\frac{5}{n\ln(n)}$.
                \item[(d)] $\sum_{n=1}^{\infty}\frac{1}{\sqrt{n^{2}+1}}$.   
            \end{enumerate}
        \end{multicols}
    \end{mdframed}

    \begin{solucion}
        \begin{enumerate}
            \item[(a)] Aplicamos el criterio del cociente. Sea
            \begin{equation*}
                a_{n}=\frac{c^{n}}{n!}.
            \end{equation*}
            Entonces
            \begin{equation*}
                \abs{\frac{a_{n+1}}{a_{n}}}=\abs{\frac{c^{n+1}}{(n+1)!}\frac{n!}{c^{n}}}=\frac{\abs{c}}{n+1}.
            \end{equation*}
            Por lo tanto,
            \begin{equation*}
                \lim_{n\to\infty}\abs{\frac{a_{n+1}}{a_{n}}}=\lim_{n\to\infty}\frac{\abs{c}}{n+1}=0\lt1.
            \end{equation*}
            Luego, por el criterio del cociente, la serie converge absolutamente para todo $c\in\R$. Además,
            \begin{equation*}
                \sum_{n=0}^{\infty}\frac{c^{n}}{n!}=e^{c}.
            \end{equation*}
            \item[(b)] Comparamos con la serie armonica. Notemos que
            \begin{equation*}
                \frac{n}{n^{2}+1}=\frac{1}{n+\frac{1}{n}}.
            \end{equation*}
            Calculamos el limite
            \begin{equation*}
                \lim_{n\to\infty}\frac{\frac{n}{n^{2}+1}}{\frac{1}{n}}=\lim_{n\to\infty}\frac{1}{n+\frac{1}{n}}.
            \end{equation*}
            Como
            \begin{equation*}
                \sum_{n=1}^{\infty}\frac{1}{n}
            \end{equation*}
            diverge, por el criterio de comparación por límite también diverge
            \begin{equation*}
                sum_{n=1}^{\infty}\frac{n}{n^{2}+1}.
            \end{equation*}
            \item[(c)] Usamos el criterio integral. Consideremos
            \begin{equation*}
                f(x)=\frac{5}{x\ln(x)},\qquad x\geq2.
            \end{equation*}
            La función es positiva, continua y decreciente para $x\geq2$. Entonces
            \begin{equation*}
                \int_{2}^{\infty}\frac{5}{x\ln(x)}dx.
            \end{equation*}
            Haciendo el cambio de variable
            \begin{equation*}
                u=\ln(x),\qquad du=\frac{dx}{x},
            \end{equation*}
            obtenemos
            \begin{equation*}
                \int_{2}^{\infty}\frac{5}{x\ln(x)}dx=\int_{\ln(2)}^{\infty}\frac{1}{u}du.
            \end{equation*}
            Pero
            \begin{equation*}
                \int_{\ln(2)}^{\infty}\frac{1}{u}du=\lim_{k\to\infty}\ln(u)\Big|_{\ln 2}^{b}=\infty.
            \end{equation*}
            Así, la serie diverge.
            \item[(d)] Comparamos con la serie armonica
            \begin{equation*}
                \lim_{n\to\infty}\frac{\frac{1}{\sqrt{n^{2}+1}}}{\frac{1}{n}}=\lim_{n\to\infty}\frac{n}{\sqrt{n^{2}+1}}=\lim_{n\to\infty}\frac{1}{\sqrt{1+\frac{1}{n^{2}}}}=1.
            \end{equation*}
            Como
            \begin{equation*}
                \sum_{n=1}^{\infty}\frac{1}{n}
            \end{equation*}
            diverge, el criterio de comparación implica que la serie tambien diverge.
        \end{enumerate}
    \end{solucion}

    \begin{mdframed}[style=mdpinkbox,frametitle={Ejercicio 5}]
        Realiza lo siguiente:
        \begin{enumerate}
            \item[(a)] Muestra que si $0.a_{1}\dots a_{n}$ es una expanción decimal finita con $a_{n}\neq0$,
            entonces
            \begin{equation*}
                0.a_{1}\dots a_{n}=0.a_{1}\dots a_{n-1}(a_{n}-1)\overline{9}.
            \end{equation*}
            \item[(b)] Muestra que si $0.a_{1}\dots a_{n}\overline{9}$ es una expansión decimal con
            $a_{n}\neq9$, entonces
            \begin{equation*}
                0.a_{1}\dots a_{n}\overline{9}=0.a_{1}\dots a_{n-1}(a_{n}+1).
            \end{equation*}
        \end{enumerate}
    \end{mdframed}

    \begin{proof}
        \begin{enumerate}
            \item[(a)] Sea $x=0.a_{1}a_{2}\dots a_{n}$, $a_{n}\neq0$. Por definición de expansión defimal infinita,
            \begin{equation*}
                x=\sum_{k=1}^{n}\frac{a_{k}}{10^{k}}.
            \end{equation*}
            Consideremos el numero
            \begin{equation*}
                y=0.a_{1}\dots a_{n-1}(a_{n}-1)\overline{9}.
            \end{equation*}
            Entonces
            \begin{equation*}
                y=\sum_{k=1}^{n-1}\frac{a_{k}}{10^{k}}+\frac{a_{n}-1}{10^{n}}+\sum_{k=n+1}^{\infty}\frac{9}{10^{k}}.
            \end{equation*}
            La ultima suma geometrica
            \begin{equation*}
                \sum_{k=n+1}^{\infty}\frac{9}{10^{k}}=\frac{9}{10^{n+1}}\sum_{j=0}^{\infty}\frac{1}{10^{j}}.
            \end{equation*}
            Como
            \begin{equation*}
                \sum_{j=0}^{\infty}\frac{1}{10^{j}}=\frac{1}{1-\frac{1}{10}}=\frac{10}{9}
            \end{equation*}
            obtenemos
            \begin{equation*}
                \sum_{k=n+1}^{\infty}\frac{9}{10^{k}}=\frac{9}{10^{n+1}}\cdot\frac{10}{9}=\frac{1}{10^{n}}.
            \end{equation*}
            sustituyendo
            \begin{equation*}
                y=\sum_{k=1}^{n-1}\frac{a_{k}}{10^{k}}+\frac{a_{n}-1}{10^{n}}+\frac{1}{10^{n}}.
            \end{equation*}
            Por lo tanto,
            \begin{equation*}
                y=\sum_{k=1}^{n-1}\frac{a_{k}}{10^{k}}+\frac{a_{n}}{10^{n}}=\sum_{k=1}^{n}\frac{a_{k}}{10^{k}}=x.
            \end{equation*}
            \item[(b)] Es inmediato del inciso (a).
        \end{enumerate}
    \end{proof}
    
    \begin{mdframed}[style=mdpinkbox,frametitle={Ejercicio 6}]
        Muestra que para todo $x\in(0,1]$ tiene una expansión decimal infinita.
    \end{mdframed}
    
    \begin{mdframed}[style=mdpinkbox,frametitle={Ejercicio 7}]
        Prueba que si $(f_{n})$ es una subsucesión de funciones que converge uniformemente en $A$
        y $B$, entonces $(f_{n})$ converge uniformemente en $A\cup B$.
    \end{mdframed}

    \begin{proof}
        Observemos que en $x\in A\cap B$ se tiene
        \begin{equation*}
            f_{n}(x)\to f(x)
        \end{equation*}
        y tambien
        \begin{equation*}
            f(x)_{n}=g(x).
        \end{equation*}
        Por unicidad del limite,
        \begin{equation*}
            f(x)=g(x).
        \end{equation*}
        Por tanto podemos definir una función
        \begin{equation*}
            h:A\cup B\to\R
        \end{equation*}
        como
        \begin{equation*}
            h(x)=
            \begin{cases}
                f(x),\quad x\in A,\\
                g(x),\quad x\in B.
            \end{cases}
        \end{equation*}
        Sea ahora $\e\gt0$. Como $f_{n}\to f$ uniformemente en $A$, existe $N_{1}\in\N$ tal que
        \begin{equation*}
            n\geq N_{1}\implies\abs{f_{n}(x)-f(x)}\lt\e\qquad\forall x\in A.
        \end{equation*}
        Así, como $f_{n}\to g$ uniformemente en $B$, existe $N_{2}\in\N$ tal que
        \begin{equation*}
            n\geq N_{2}\implies \abs{f_{n}(x)-g(x)}\lt\e\qquad x\in B.
        \end{equation*}
        Tomemos $\max\{N_{1},N_{2}\}$. Entonces, para todo $n\gt N$ y todo $x\in A\cup B$, ocurren dos posibilidades. Si
        $x\in A$, entonces
        \begin{equation*}
            \abs{f_{n}(x)-h(x)}=\abs{f_{n}-f(x)}\lt\e.
        \end{equation*}
        Si $x\in B$, entonces
        \begin{equation*}
            \abs{f_{n}(x)-h(x)}=\abs{f_{n}-g(x)}\lt\e.
        \end{equation*}
        Por ende
        \begin{equation*}
            \abs{f_{n}(x)-h(x)}\lt\qquad\forall x\in A\cup B,
        \end{equation*}
        Para todo $n\geq N$. Por lo tanto, $(f_{n})$ converge uniformemente en $A\cup B$.
    \end{proof}

\end{document}

